\chapter{Konzept}
\section{Grundlagen der Mathematik}

Die Mathematik ist eines der mächtigsten Werkzeuge in \LaTeX{}. Wir können sowohl inline-Mathematik wie $E = mc^2$ als auch Display-Mathematik verwenden.

\begin{figure}[h]
\centering
\fbox{\begin{minipage}{8cm}
Dies ist ein Platzhalter für eine Grafik.\\
Mathematische Darstellung von Konzepten
\end{minipage}}
\caption{Visualisierung mathematischer Konzepte}
\label{fig:math_visualization}
\end{figure}

\subsection{Quadratische Gleichungen}

Die allgemeine Form einer quadratischen Gleichung ist:
\[
ax^2 + bx + c = 0
\]

Die Lösungsformel lautet:
\[
x = \frac{-b \pm \sqrt{b^2 - 4ac}}{2a}
\]

\subsection{Differenzial- und Integralrechnung}

Die Ableitung einer Funktion wird definiert als:
\[
f'(x) = \lim_{h \to 0} \frac{f(x+h) - f(x)}{h}
\]

Das bestimmte Integral:
\[
\int_a^b f(x) \, dx
\]

\section{Physik}

\subsection{Klassische Mechanik}

Newtons zweites Gesetz besagt:
\[
\vec{F} = m \vec{a}
\]

Die kinetische Energie ist:
\[
E_k = \frac{1}{2}mv^2
\]

\begin{figure}[h]
\centering
\fbox{\begin{minipage}{8cm}
Grafische Darstellung von Kräften\\
und Bewegungsvektoren
\end{minipage}}
\caption{Kräfte und Beschleunigungen in der klassischen Mechanik}
\label{fig:classical_mechanics}
\end{figure}
